\label{sec:conclusion}

M.1 is the first community character paper in the M-track series, establishing
economic character---commercial intensity, industrial fraction, and jobs per
resident---as an operationalizable signal derived from the Census Bureau's
LODES Workplace Area Characteristics data.
The signal fits the M-track framework (M.0) as a cosine-similarity plug-in
for the edge weight modifier, and composes with all existing \textsc{bisect}
weight modes and structure algorithms.

The empirical results demonstrate clear state-specificity in the economic
character signal's effectiveness.
North Carolina's concentrated employment clusters (Research Triangle Park,
Charlotte commercial corridor) act as strong economic community boundaries
that the algorithm can exploit, yielding a 14.3\,pp proportionality gap
improvement.
Wisconsin's diffuse economic geography produces no effect, as the mathematical
framework predicts.
Texas's industrial zones are embedded in Republican-leaning suburban areas,
producing a slight reversal.

These results establish the M-track hypothesis that community character
signals should be deployed selectively, where the community geography is
concentrated and identifiable.
A community-of-interest signal that has no measurable effect is not a
communities-of-interest tool; it is noise.
M.1 passes the Wisconsin null test, which is as important as the NC
positive result.

\paragraph{Relationship to the M-track series.}
M.1 establishes the empirical pattern that will guide the development of
M.2 through M.7.
The key design question for each subsequent paper is:
does this state have a concentration of the relevant community type
(land-use clusters for M.2, housing character for M.3, commuting sheds for M.4,
topographic zones for M.5, administrative zone co-memberships for M.6,
transit networks for M.7) that would make the signal mechanistically effective?
The M-track framework answer is: run the null test first.
If the signal produces no effect on Wisconsin, it is not useful for Wisconsin
regardless of how legally well-grounded it is.

\paragraph{Phase 2 work.}
The M-track primary metric $\overline{\sigma}_{\mathrm{JPR}}$ (within-district
jobs-per-resident standard deviation) requires per-district tract-level output
not produced in the current run.
Phase~2 will compute this metric for all three test states and extend the
evaluation to additional states with concentrated employment clusters:
Massachusetts (biomedical corridor), California (Silicon Valley, port
complexes), Pennsylvania (Pittsburgh industrial legacy).
Phase~2 will also provide multi-seed variance estimates for the WI and TX
single-seed results.

\paragraph{Integration with M.8 composite.}
The economic character signal from M.1 will contribute one dimension to
the M.8 composite community character index (not yet published).
The M.8 paper will provide an expert witness template that combines M.1--M.7
signals with tunable weights and maps the composite score to specific
communities of interest recognized in redistricting proceedings.
M.1 is expected to be a high-weight component in any state with identifiable
employment concentration clusters.
